
%  The lines beginning with "%" are comments. 
%  Theses comments explain how to use this template for
%  preparing the Seminar report. 
%
%  The first non-comment line below contain the matter that
%  go into the preamble of the input file. These are collected in a 
%  file by name "preamble.tex". This file can be seen in the 
%  folder "SeminarReportTemplate". 
%
%  Do not make any changes in the first non-comment line below.
%
\documentclass[11pt, twoside]{book}
\let\cleardoublepage\clearpage
\usepackage{graphicx, fancyhdr, amsmath,times, url}
\setlength{\oddsidemargin}{3cm}
\newcommand{\VAtitle}[1]{\def\vtitle{Bitcoin: A Cryptocurrency}}
\newcommand{\VAauthor}[1]{\def\vauthor{Savitha K J}}% Write full name without, repeat without, Mr. or Ms.
\newcommand{\VAadmissionyear}[1]{\def\vadmissionyear{2010}}
\newcommand{\VAregisternumber}[1]{\def\vregisternumber{VEAKECS049}}% Write University Examination Register Number.
\newcommand{\VAguide}[1]{\def\vguide{Ms. Geethu P C M.Tech, MISTE}}% Write full name with Mr. or Ms. or Dr. or Prof.
\newcommand{\VAhod}[1]{\def\vhod{Col MP Induchudan,  M.E., FITE}}% Write full name with Mr. or Ms. or Dr. or Prof. 
%
%
%  No changes in the next non-comment line also.
%
\begin{document}
%
%***************************************************
%
%  In the next line, replace "Seminar Title" by the title of your seminar.
%
\VAtitle{Seminar Title}%
%
%  In the next line, replace "Student" by your name. 
%  Write full name without, repeat without, Mr. or Ms.
%
\VAauthor{Student}%
%
%  In the next line, replace "2008" by the year of your admission 
%  to the College.
%
\VAadmissionyear{2010}%
%
%  In the next line, replace "VExxMCAnnn" by your 
%  University Examination Register Number.
%
\VAregisternumber{VExxxCSnnn}% 
%
%  In the next line, replace "Guide" by the full name of your guide or supervisor 
%  with Mr. or Ms. or Dr. or Prof.
%
\VAguide{Guide}% 
%
%  In the next line, replace "hod" by the full name of Head of Department 
%  with Mr. or Ms. or Dr. or Prof.
\VAhod{hod}
%
%  Printing the first few pages: title, certificate, acknowledgement, etc.
%  No changes in the next non-comment line.
%
%
\makeatletter
\renewcommand{\chapter}{\if@openright\cleardoublepage\else\clearpage\fi
                    %\thispagestyle{fancy}%
                    \global\@topnum\z@
                    \@afterindentfalse
                    \secdef\@chapter\@schapter}%
%
\makeatother
%
\pagestyle{empty}
%\frontmatter
%
%******************************************************
%
\begin{titlepage}
\newcommand{\HRule}{\rule{\linewidth}{0.5mm}}
\begin{center}
% Upper part of the page
% Title
%\quad\\[0.3cm]
{\large \sffamily CS 09 805 (P) : Seminar}\\[0.2cm]
\HRule \\[0.5cm]
{ \huge\sffamily \bfseries \vtitle}\\[0.2cm]
\HRule \\[5.20cm]
% Author
{\Large \sffamily\bfseries  \vauthor}\\
{\large \sffamily Reg. No. \vregisternumber\\ 
S8 CSE (\vadmissionyear\ Admissions)} \\[4cm] 
% Bottom of the page
%\vfill
\includegraphics[width=0.40\textwidth]{VidyaLogo}\\[0.2cm]
{\Large \sffamily\bfseries Department of Computer Science and Engineering}\\ {\large\sffamily Vidya Academy of Science \& Technology\\ \normalsize Thalakkottukara, Thrissur - 680 501}\\
({\tt http://www.vidyaacdemy.ac.in})\\
\end{center}
\end{titlepage}
%
%******************************************************
%

%%
%******************************************************
%
\begin{titlepage}
\newcommand{\HRule}{\rule{\linewidth}{0.5mm}}
\begin{center}
% Upper part of the page
% Title
%\quad\\[0.3cm]
{\large \sffamily CS 09 805 (P) : Seminar}\\[1cm]
%\HRule \\[0.5cm]
{ \huge\sffamily \bfseries \vtitle}\\[5.9cm]
% Author
{\Large \sffamily \bfseries \vauthor}\\
{\large \sf  Reg. No. \vregisternumber\\ 
S8 CSE (\vadmissionyear\  Admissions)} \\[4cm] 
% Bottom of the page
%\vfill
\includegraphics[width=0.40\textwidth]{VidyaLogo}\\[0.2cm]
{\Large \sffamily \bfseries Department of Computer Science and Engineering}\\ {\large \sffamily Vidya Academy of Science \& Technology\\ \normalsize Thalakkottukara, Thrissur - 680 501}\\
({ \bf \tt http://www.vidyaacademy.ac.in})\\[1cm]
\end{center}
\end{titlepage}
%
%******************************************************
%
\clearpage
\pagenumbering{roman}
%
\addcontentsline{toc}{chapter}{\quad Certificate}
%
\begin{center}
\hrule
\quad\\[0.75cm]
{\Large \bf Department of Computer Science and Technology}\\
{\large \bf Vidya Academy of Science \& Technology}\\
{\normalsize \bf Thalakkottukara, Thrissur - 680 501\\
({\tt http://www.vidyaacademy.ac.in})}\\[0.5cm]
\includegraphics{VidyaLogo}\\[1.25cm]
%
{\Large \bf CERTIFICATE}\\[1cm]
%
\end{center}
\index{certificate}
\index{University of Calicut}
This is to certify that the report titled 
{\bf \vtitle} is a bona-fide record of the 
Seminar presented by {\bf \vauthor\  
(Reg. No. \vregisternumber)}    of the Eighth 
Semester B.Tech CSE (\vadmissionyear\ admissions) 
of Vidya Academy of Science \& Technology, 
Thrissur - 680 501 in partial fulfilment of the 
requirement for the award of the Bachelor of Technology in Computer Science and Engineering from the University of 
Calicut.\\[0.1cm]
%
\begin{center}
\begin{tabular}{llll}
\multicolumn{2}{l}{\bf Seminar Guide}    
&\multicolumn{2}{l}{\bf Head of Department}   \\
   & & \\
Name: &&  Name: \\
\vguide\ && \vhod\ \\ \\
Signature: \quad  ................................
\qquad \qquad && Signature: \quad 
................................ \\
& & & \\
Date: \hspace{10mm} \today & & Date: \hspace{10mm} \today \\
\end{tabular}
\\[1cm]
{\small (Seal of Department of Computer Science and Engineering)}
\quad \\[0.75cm]
\hrule
\end{center} 
%
\vspace*{\fill}
%
\clearpage
\pagestyle{fancy}
\lhead{{\small \vtitle}} \rhead{\thepage}
\lfoot{{\scriptsize Department of Computer Science \& Engineering, 
Vidya Academy of Science \& Technology, Thrissur}}
\rfoot{\includegraphics[width=0.5cm]{VidyaLogo}}
\cfoot{ }
%
\renewcommand{\headrulewidth}{1pt}
\renewcommand{\footrulewidth}{0.5pt}
%
%\clearpage
\thispagestyle{fancy}
%
\addcontentsline{toc}{chapter}{\quad Acknowledgement}
%
\chapter*{Acknowledgement}
%
\index{acknowledgment}
I wish to record my indebtedness and thankfulness 
to all who helped me prepare this Seminar Report\index{seminar report} titled 
\vtitle\  and present it in a satisfactory way. 

My sincere thanks to Dr.Sudha Balagopalan, our Principal, for providing us all the necessary facilities.
I am also thankful to \vhod, 
the Head of Department of Computer Science and Engineering 
for encouragement and for giving valuable 
support.
I am especially thankful to my seminar 
guide \vguide\  
for giving me valuable suggestions and 
critical inputs in the preparation of this paper. 

My classmates have always been 
helpful and I am grateful to them for 
patiently listening to  my presentation of the seminar and interacting with me through their valuable questions. 


%
\clearpage
%

%
%
%****************************************************
%
% No changes in the next two non-comment lines.
%
\chapter*{Abstract}
\addcontentsline{toc}{chapter}{\quad Abstract}
%
%  Here enter the abstract of the paper.
Bitcoin is the world's first decentralized digital currency, allowing the easy 
storage and transfer of cryptographic tokens, using a peer-to-peer network to carry
information, hashing as a synchronization signal to prevent double-spending, 
and a powerful scripting system to determine ownership of the tokens. There is a 
growing technology and business infrastructure supporting it.


A purely peer-to-peer version of electronic cash would allow online
payments to be sent directly from one party to another without going through a
financial institution. Digital signatures provide part of the solution, but the main
benefits are lost if a trusted third party is still required to prevent double-spending.

The network timestamps transactions by hashing them into an ongoing chain of
hash-based proof-of-work, forming a record that cannot be changed without 
redoing the proof-of-work. The longest chain not only serves as proof of the 
sequence of events witnessed, but proof that it came from the largest pool of CPU 
power. As long as a majority of CPU power is controlled by nodes that are not 
cooperating to attack the network, they'll generate the longest chain and outpace 
attackers. The network itself requires minimal structure. Messages are broadcast on 
a best effort basis, and nodes can leave and rejoin the network at will, accepting the 
longest proof-of-work chain as proof of what happened while they were gone.

To send money, you broadcast to the network that the amount on your account 
should go down, and the amount on a receiver’s account up. Nodes, or computers,
in the Bitcoin network apply that transaction to their copy of the ledger, and then 
pass on the transaction to other nodes. This, with some math-based security, is 
really all there is--a system that lets a group of computers maintain a ledger.

%
%***************************************************
%
%  The next line creates a Table of Contents. 
%  Do not delete this. There must be a Table of Contents.
%
\tableofcontents
%
%  The next line creates a page containing a List of Tables.
%  Delete it if there are no tables in your paper. 
%
\listoftables
%
%  The next line creates a page containing a List of Figures.
%  Delete it if there are no figures in your paper.
%
\listoffigures
%
%***************************************************
%
\mainmatter
%
%  The main contents of the paper begin here. 
%  In the following replace "Abcde Xyz" with 
%  the title of the first chapter of your paper.
%
\chapter{Introduction}
%
%  The text given below is only some filler text. Delete 
%  all those and replace them with the contents of your % 
%  paper. It is not exactly gibberish. (The filler text is a 
%  passage in Latin written by Cicero.)
%
Until Bitcoin’s invention in 2008 by the unidentified programmer known as Satoshi Nakamoto, online 
transactions always required a trusted third-party intermediary. For example, if Alice wanted to send \$100 
to Bob over the Internet, she would have had to rely on a third-party service like PayPal or MasterCard. 
Intermediaries like PayPal keep a ledger of account holders’ balances. When Alice sends Bob \$100, 
PayPal deducts the amount from her account and adds it to Bob’s account. Without such intermediaries, 
digital money could be spent twice. Imagine there are no intermediaries with ledgers, and digital cash 
is simply a computer file, just as digital documents are computer files. Alice could send \$100 to Bob 
by attaching a money file to a message. But just as with email, sending an attachment does not remove 
it from one’s computer. Alice would retain a copy of the money file after she had sent it. She could 
then easily send the same \$100 to Charlie. In computer science, this is known as the “double-spending” 
problem and until Bitcoin it could only be solved by employing a ledger-keeping trusted third party. Bit 
coin’s invention is revolutionary because for the first time the double-spending problem can be solved 
without the need for a third party. Bitcoin does this by distributing the necessary ledger among all the 
users of the system via a peer-to-peer network. Every transaction that occurs in the Bitcoin economy is 
registered in a public, distributed ledger, which is called the block chain. New transactions are checked 
against the block chain to ensure that the same bit coins haven’t been previously spent, thus eliminating 
the double-spending problem. The global peer-to-peer network, composed of thousands of users, takes the 
place of an intermediary; Alice and Bob can transact without PayPal. One thing to note right away is that 
transactions on the Bitcoin network are not denominated in dollars or Euros or yens they are on PayPal, 
but are instead denominated in Bitcoin. This makes it a virtual currency in addition to a decentralized 
payments network. The value of the currency is not derived from gold or government fiat, but from the 
value that people assign to it. The dollar value of a Bitcoin is determined on an open market, just as is the 
exchange rate between different world currencies.
%
%  Replace "Fghi Jklm" with the title of 
%  the first section of the first chapter of your Report.
%

\chapter{Related Works}
\textbf{Electronic funds transfer (EFT)} is the electronic exchange, transfer of money from one account to
another, either within a single financial institution or across multiple institutions, through computer-based 
systems.
It consists of different concepts;

\section{Fghi Jklm}
\subsection{yo mayo}
%
Sed ut perspiciatis, unde omnis iste natus error sit 
voluptatem accusantium doloremque laudantium, totam 
rem aperiam eaque ipsa, quae ab illo inventore 
veritatis et quasi architecto beatae vitae dicta 
sunt, explicabo. Nemo enim ipsam voluptatem, quia 
voluptas sit, aspernatur aut odit aut fugit, sed quia 
consequuntur magni dolores eos, qui ratione voluptatem 
sequi nesciunt, neque porro quisquam est, qui dolorem 
ipsum, quia dolor sit amet, consectetur, adipisci[ng] 
velit, sed quia non numquam [do] eius modi tempora 
inci[di]dunt, ut labore et dolore magnam aliquam 
quaerat voluptatem. Ut enim ad minima veniam, quis 
nostrum exercitationem ullam corporis suscipit laboriosam, 
nisi ut aliquid ex ea commodi consequatur? Quis autem vel 
eum iure reprehenderit, qui in ea voluptate velit esse, 
quam nihil molestiae consequatur, vel illum, qui dolorem 
eum fugiat, quo voluptas nulla pariatur?
%
\section{Opqrs Tuvw}
%
At vero eos et accusamus et iusto odio dignissimos 
ducimus, qui blanditiis praesentium voluptatum 
deleniti atque corrupti, quos dolores et quas 
molestias excepturi sint, obcaecati cupiditate 
non provident, similique sunt in culpa, qui officia 
deserunt mollitia animi, id est laborum et dolorum 
fuga. Et harum quidem rerum facilis est et expedita 
distinctio. Nam libero tempore, cum soluta nobis 
est eligendi optio, cumque nihil impedit, quo 
minus id, quod maxime placeat, facere possimus, 
omnis voluptas assumenda est, omnis dolor repellendus. 
Temporibus autem quibusdam et aut officiis debitis 
aut rerum necessitatibus saepe eveniet, ut et 
oluptates repudiandae sint et molestiae non 
recusandae. Itaque earum rerum hic tenetur a 
sapiente delectus, ut aut reiciendis voluptatibus 
maiores alias consequatur aut perferendis doloribus 
asperiores repellat...
%
% Continue adding new chapters, sections and 
% subsections.
%
% ***   Creating Bibliography   ***
%
% Do not delete the following two lines
%
\clearpage
\addcontentsline{toc}{chapter}{\quad Bibliography}
%
% The markups for creating the bibliography begins here.
% One sample bibliography item is included. 
% Add additional items if there are any.
%
\begin{thebibliography}{99}
%
\bibitem{First}
{\em First Lessons in \LaTeX}, by Dr. V N Krishnachandran, 
Vidya Academy of Science \& Technology, 
Thrissur - 680 501, 2011.
%
\end{thebibliography}
%
% Printing the last page
%
%
\newpage
\thispagestyle{empty}
\vspace*{\fill}
\begin{flushright}
\includegraphics{VidyaLogo}\\[0.5cm]
{\Large \bf \sf  Department of Computer Science and Engineering}\\
{\sf Vidya Academy of Science \& Technology\\
Thalakkottukara, Thrissur - 680 501\\
({\tt http://www.vidyaacademy.ac.in})}
\end{flushright}
%
%
%   ***   The end   ***
%
\end{document}
